% =====================================================================
%  Bölüm 10 — Önermeler Mantığının Sözdizimi ve Anlambilimi
% =====================================================================
\chapter{Sözdizimi ve Anlambilim}\label{bol:sozdizim}

\begin{bolumhedefleri}
\begin{itemize}
  \item \emph{Nesne dili} ile \emph{üst dil} ayrımını ve Kısım~III'ün
        neden bu ayrıma dayandığını,
  \item Önermeler mantığının formüllerini \Lean{}'de bir tümevarımsal
        tip olarak tanımlamayı,
  \item Formüller üzerinde yapısal özyinelemeyi: boyut, atomlar,
        yerine koyma,
  \item Değerlemeleri ve Tarski'nin doğruluk tanımını,
  \item Totoloji, sağlanabilirlik ve mantıksal sonuç ($\models$)
        kavramlarını,
  \item Doğruluk tablolarına dayanan bir karar yöntemini ve sınırlarını,
  \item Doğruluk tablosu anlambiliminin \emph{klasik} olduğunu.
\end{itemize}
\end{bolumhedefleri}

% ---------------------------------------------------------------------
\section{Kısım III'e giriş: mantık hakkında mantık}\label{sec:kisim3-giris}

Şimdiye kadar mantığı \emph{kullandık}: \Lean{}'de önermeler yazdık ve
onları ispatladık. Kısım~III'te mantığın kendisini bir matematiksel
\emph{nesne} olarak inceleyeceğiz. Soracağımız sorular, mantık
\emph{içinde} değil, mantık \emph{hakkında} sorulardır:

\begin{itemize}
  \item İspatlayabildiğimiz her şey doğru mudur? (\emph{geçerlilik},
        \cref{bol:gecerlilik})
  \item Doğru olan her şeyi ispatlayabilir miyiz? (\emph{tamlık},
        \cref{bol:tamlik})
  \item Her ispat dolaysız bir ispata dönüştürülebilir mi?
        (\emph{kesme eleme}, \cref{bol:ardisik})
\end{itemize}

\noindent Bu sorulara verilen cevaplara \emph{metamantık} denir.
\Cref{fig:buyuk-resim}, Kısım~III'ün büyük resmini gösteriyor.

\begin{figure}[htbp]
\centering
\begin{tikzpicture}[
    kutu/.style={draw=#1, fill=#1!8, very thick, rounded corners=6pt,
                 text width=4.2cm, align=center, inner sep=8pt},
    ok/.style={-{Stealth[length=3mm]}, very thick, #1},
  ]
  \node[kutu=anarenk] (s) at (0,0)
    {{\baslikfont\bfseries\color{anarenk} SÖZDİZİMİ}\\[3pt]
     $\Gamma \vdash A$\\[2pt]
     {\small ``$A$, $\Gamma$'dan \emph{türetilebilir}''}\\[2pt]
     {\scriptsize çıkarım kuralları, ispat ağaçları}};
  \node[kutu=ikincil] (a) at (10,0)
    {{\baslikfont\bfseries\color{ikincil} ANLAMBİLİM}\\[3pt]
     $\Gamma \models A$\\[2pt]
     {\small ``$A$, $\Gamma$'nın \emph{mantıksal sonucudur}''}\\[2pt]
     {\scriptsize değerlemeler, doğruluk tabloları}};
  \draw[ok=teoremframe] ([yshift=0.55cm]s.east) -- node[above, font=\baslikfont\small\bfseries, text=teoremframe]{Geçerlilik (Bölüm 11)} ([yshift=0.55cm]a.west);
  \draw[ok=vurgu] ([yshift=-0.55cm]a.west) -- node[below, font=\baslikfont\small\bfseries, text=vurgu]{Tamlık (Bölüm 12)} ([yshift=-0.55cm]s.east);
  \node[font=\baslikfont\scriptsize, text=black!60] at (5,1.3) {ispatlanabilen her şey doğrudur};
  \node[font=\baslikfont\scriptsize, text=black!60] at (5,-1.3) {doğru olan her şey ispatlanabilir};
  \node[font=\baslikfont\small\bfseries, text=koyu] at (5,-2.4) {Bu bölüm: iki tarafın da tanımları};
\end{tikzpicture}
\caption{Kısım III'ün büyük resmi. Sözdizimi ile anlambilim, mantığın
iki yüzüdür; geçerlilik ve tamlık teoremleri bu iki yüzü birbirine bağlar.}
\label{fig:buyuk-resim}
\end{figure}

% ---------------------------------------------------------------------
\section{Nesne dili ve üst dil}\label{sec:nesne-ust}

Mantık hakkında konuşurken iki dili birbirinden ayırmak zorundayız:

\begin{tanim}{Nesne dili ve üst dil}{nesne-ust}
\textbf{Nesne dili}, incelediğimiz mantık sistemidir: formülleri,
bağlaçları, çıkarım kuralları. \textbf{Üst dil} ise nesne dili hakkında
konuşmak ve onun özelliklerini ispatlamak için kullandığımız dildir.
\end{tanim}

\noindent Kâğıt üzerindeki mantık kitaplarında iki dil de aynı sembollerle
yazıldığı için karışıklık kolayca ortaya çıkar. ``$A \to B$ ispatlanabilirse
ve $A$ ispatlanabilirse $B$ ispatlanabilir'' cümlesindeki ilk ``$\to$''
nesne dilinin bir bağlacıdır; ``ise''ler ise üst dilin parçasıdır.

\Lean{}'de bu ayrım kendiliğinden ortaya çıkar (\cref{fig:nesne-ust}).
Üst dil \Lean{}'in kendisidir: \lc{Prop}, \lc{∀}, \lc{→}. Nesne dili ise
\Lean{}'de tanımlayacağımız bir \emph{veri tipidir}: \lc{Formul}. Bir
formül, bir liste ya da bir ağaç gibi sıradan bir veridir; onu
hesaplayabilir, dönüştürebilir ve onun hakkında teoremler
ispatlayabiliriz.

\begin{figure}[htbp]
\centering
\begin{tikzpicture}[
    katman/.style={draw=#1, fill=#1!6, very thick, rounded corners=8pt},
  ]
  \draw[katman=anarenk] (0,0) rectangle (11.4,4.2);
  \node[anchor=north west, font=\baslikfont\bfseries, text=anarenk] at (0.2,4.1) {ÜST DİL: Lean};
  \node[anchor=north west, font=\small, text=black!70, align=left] at (0.2,3.55)
       {\lc{Prop}, \lc{∀}, \lc{∃}, \lc{→}, \lc{theorem}\\[1pt] \lc{Totoloji A}, \lc{Sonuc Γ A}};
  \draw[katman=ikincil] (5.2,0.4) rectangle (11.0,3.3);
  \node[anchor=north west, font=\baslikfont\bfseries, text=ikincil] at (5.4,3.2) {NESNE DİLİ: \lc{Formul}};
  \node[anchor=north west, font=\small, align=left] at (5.4,2.6)
       {\lc{atom 0}, \lc{bot}\\[2pt]
        $p \oise q$, \quad $p \ove q$, \quad $p \oveya {\sim}p$\\[2pt]
        {\scriptsize sıradan veri: hesaplanır, dönüştürülür}};
  \draw[-{Stealth[length=3mm]}, very thick, vurgu] (3.4,1.2) -- node[above, font=\baslikfont\scriptsize\bfseries, text=vurgu, align=center]{hakkında\\teorem} (5.1,1.2);
  \node[font=\leanfont\scriptsize, align=left, anchor=west] at (0.2,1.2)
       {theorem peirce\_totoloji :\\\quad Totoloji peirce};
\end{tikzpicture}
\caption{Üst dil (\Lean{}) ile nesne dili (\lc{Formul} tipi). Nesne
dilinin formülleri, üst dilde birer veridir.}
\label{fig:nesne-ust}
\end{figure}

% ---------------------------------------------------------------------
\section{Formüllerin sözdizimi}\label{sec:formul-sozdizim}

Önermeler mantığının formülleri, kâğıt üzerinde genellikle şöyle bir
\emph{dilbilgisiyle} tanımlanır:
\[
  A, B \;::=\; p_n \;\mid\; \bot \;\mid\; A \wedge B \;\mid\; A \vee B \;\mid\; A \to B
\]
``Bir formül ya bir atomdur, ya $\bot$'dur, ya da iki formülün
bağlaçla birleştirilmesidir.'' Bu tanımın \Lean{}'deki karşılığı,
harfi harfine bir tümevarımsal tiptir:

\leandosya[Formüller]{Bolum10}{formul_tanim}

\noindent Atomları doğal sayılarla adlandırıyoruz: $p_0, p_1, p_2, \ldots$
Değillemeyi ayrı bir kurucu olarak eklemiyoruz; Bölüm~3'teki gibi
$\neg A := A \to \bot$ olarak tanımlıyoruz. Okunabilirlik için nesne
dilinin bağlaçlarına \Lean{}'in kendi bağlaçlarından \emph{farklı}
semboller veriyoruz:

\leandosya[Gösterim ve tanımlanmış bağlaçlar]{Bolum10}{gosterim}

\begin{center}
\renewcommand{\arraystretch}{1.4}
\rowcolors{2}{tablosatir}{white}
\begin{tabular}{>{\centering\arraybackslash}p{2.4cm}
                >{\centering\arraybackslash}p{2.6cm}
                >{\raggedright\arraybackslash}p{6.4cm}}
\rowcolor{tabloust}
\tablobaslik{Nesne dili} & \tablobaslik{Üst dil (Lean)} & \tablobaslik{Fark} \\
$A \ove B$     & \lc{P ∧ Q}  & biri bir \lc{Formul} (veri), öteki bir \lc{Prop} \\
$A \oveya B$   & \lc{P ∨ Q}  & aynı \\
$A \oise B$    & \lc{P → Q}  & aynı \\
${\sim}A$      & \lc{¬P}     & aynı \\
\end{tabular}
\end{center}

\begin{dikkat}[Dikkat: Formüller önerme değildir]
$p \ove q$ yazdığımızda bir \emph{önerme} değil, bir \emph{veri} (bir
sözdizim ağacı) kurarız. Bu ağaç kendi başına ne doğrudur ne yanlış;
\lc{Formul} tipinin bir elemanıdır, tıpkı \lc{[1, 2, 3]}'ün bir liste
olması gibi. Bir formülün doğru olup olmadığını sorabilmek için önce ona
bir \emph{anlam} vermemiz gerekir (\cref{sec:anlambilim}).
\end{dikkat}

\subsection{Sözdizim ağaçları}

Bir formül, aslında bir ağaçtır. Parantezler ve öncelik kuralları
yalnızca bu ağacı tek satıra yazmanın yollarıdır
(\cref{fig:sozdizim-agaci}).

\leandosya[Örnek formüller]{Bolum10}{ornek_formul}

\begin{figure}[htbp]
\centering
\begin{tikzpicture}[
    baglac/.style={circle, draw=leansym, fill=leansym!10, thick,
                   minimum size=0.8cm, font=\small},
    atom/.style={rectangle, rounded corners=3pt, draw=leanhyp, fill=leanhyp!10,
                 thick, minimum size=0.6cm, font=\small},
    level distance=1.05cm,
    level 1/.style={sibling distance=3.0cm},
    level 2/.style={sibling distance=1.9cm},
    level 3/.style={sibling distance=1.2cm},
    edge from parent/.style={draw, black!45, thick},
  ]
  \node[baglac] {$\Rightarrow$}
    child {node[baglac] {$\Rightarrow$}
      child {node[baglac] {$\Rightarrow$}
        child {node[atom] {$p$}}
        child {node[atom] {$q$}}}
      child {node[atom] {$p$}}}
    child {node[atom] {$p$}};
  \node[font=\small, anchor=west, align=left] at (2.8,-1.2)
       {$((p \oise q) \oise p) \oise p$\\[3pt]
        {\scriptsize Peirce yasası: 3 bağlaç, 4 atom geçişi}};
\end{tikzpicture}
\caption{Peirce yasasının sözdizim ağacı. Her iç düğüm bir bağlaç,
her yaprak bir atomdur.}
\label{fig:sozdizim-agaci}
\end{figure}

\subsection{Formüller üzerinde özyineleme}

\lc{Formul} bir tümevarımsal tip olduğu için, Bölüm~9'daki ağaçlar
gibi, üzerinde fonksiyonlar özyinelemeyle tanımlanır ve özellikler
yapısal tümevarımla ispatlanır. Her kurucu için bir durum vardır;
ikili bağlaçlarda iki tümevarım hipotezi gelir.

\leandosya[Formüller üzerinde özyineleme]{Bolum10}{formul_ozyineleme}

% ---------------------------------------------------------------------
\section{Anlambilim: doğruluk}\label{sec:anlambilim}

Bir formüle anlam vermek, atomlarının doğru mu yanlış mı olduğunu
söylemekle başlar.

\begin{tanim}{Değerleme ve doğruluk değeri}{degerleme}
Bir \textbf{değerleme}, her atoma bir doğruluk değeri veren bir
fonksiyondur: $v : \mathbb{N} \to \{\mathsf{true}, \mathsf{false}\}$.
Bir $A$ formülünün $v$ altındaki \textbf{doğruluk değeri} $\lBrack A
\rBrack_v$, formülün yapısı üzerinde özyinelemeyle tanımlanır:
\begin{align*}
\lBrack p_n \rBrack_v &= v(n) &
\lBrack A \ove B \rBrack_v &= \lBrack A \rBrack_v \;\&\&\; \lBrack B \rBrack_v \\
\lBrack \bot \rBrack_v &= \mathsf{false} &
\lBrack A \oveya B \rBrack_v &= \lBrack A \rBrack_v \;\|\; \lBrack B \rBrack_v \\
& & \lBrack A \oise B \rBrack_v &= \mathord{!}\lBrack A \rBrack_v \;\|\; \lBrack B \rBrack_v
\end{align*}
\end{tanim}

\leandosya[Tarski'nin doğruluk tanımı]{Bolum10}{deger}

\noindent Bu tanımın ilginç bir yanı vardır: Nesne dilindeki ``ve''nin
anlamını üst dildeki ``ve'' (\lc{\&\&}) ile veriyoruz. Bu bir kısır döngü
değildir; çünkü üst dilin ``ve''sini zaten anlıyoruz ve onu, henüz anlamı
olmayan bir sembole ($\ove$) anlam vermek için kullanıyoruz.

\begin{tarihce}[Tarihçe: Tarski ve doğruluk]
Polonyalı mantıkçı Alfred Tarski 1933'te yayımladığı ünlü çalışmasında,
biçimsel bir dilde ``doğruluk'' kavramının nasıl tanımlanabileceğini
gösterdi \citep{tarski1933}. Tarski'nin anahtar fikri, nesne dili ile
üst dili kesin biçimde ayırmaktı: Bir dilin doğruluk kavramı, o dilin
\emph{kendi içinde} değil, daha güçlü bir üst dilde tanımlanmalıdır. Bu
kısıtlama, yalancı paradoksunu (``Bu cümle yanlıştır'') engeller.
Tarski'nin \emph{doğruluğun tanımlanamazlığı teoremi}, Gödel'in eksiklik
teoremiyle birlikte \ciltiki{4}'ün konusudur.
\end{tarihce}

\leandosya[Değillemenin değeri]{Bolum10}{deger_degil}

\subsection{Doğruluk tabloları}

Bir formülün doğruluk tablosu, atomlarına verilebilecek bütün değerler
için \lc{deger}'in sonucunu listeler. \Cref{fig:peirce-tablo}, Peirce
yasasının tablosunu gösteriyor.

\begin{figure}[htbp]
\centering
\begin{tikzpicture}
  \matrix[matrix of nodes, nodes={minimum width=1.9cm, minimum height=0.68cm,
          anchor=center, font=\small}, column sep=0pt, row sep=0pt,
          row 1/.style={nodes={fill=tabloust, text=white, font=\baslikfont\bfseries\scriptsize}}] (t) {
    $p$ & $q$ & $p \oise q$ & $(p \oise q) \oise p$ & Peirce \\
    1 & 1 & 1 & 1 & |[fill=bittibg, text=bittiframe]| \textbf{1} \\
    1 & 0 & 0 & 1 & |[fill=bittibg, text=bittiframe]| \textbf{1} \\
    0 & 1 & 1 & 0 & |[fill=bittibg, text=bittiframe]| \textbf{1} \\
    0 & 0 & 1 & 0 & |[fill=bittibg, text=bittiframe]| \textbf{1} \\
  };
  \draw[black!25] (t.north west) rectangle (t.south east);
  \draw[black!25] (t-1-3.north west) -- (t-5-3.south west);
\end{tikzpicture}
\caption{Peirce yasası $((p \oise q) \oise p) \oise p$'nin doğruluk
tablosu. Son sütun her satırda 1'dir: Formül bir totolojidir.}
\label{fig:peirce-tablo}
\end{figure}

% ---------------------------------------------------------------------
\section{Totoloji ve mantıksal sonuç}\label{sec:totoloji}

\begin{tanim}{Totoloji, sağlanabilirlik, mantıksal sonuç}{totoloji}
\begin{itemize}
  \item $A$ bir \textbf{totolojidir} ($\models A$): her $v$ için
        $\lBrack A \rBrack_v = \mathsf{true}$.
  \item $A$ \textbf{sağlanabilirdir}: en az bir $v$ için
        $\lBrack A \rBrack_v = \mathsf{true}$.
  \item $A$, $\Gamma$'nın \textbf{mantıksal sonucudur} ($\Gamma \models A$):
        $\Gamma$'daki bütün formülleri doğru yapan her $v$, $A$'yı da
        doğru yapar.
\end{itemize}
\end{tanim}

\leandosya[Tanımlar]{Bolum10}{totoloji_tanim}

\noindent Bu tanımlar \emph{üst dilde} birer önermedir: \lc{Totoloji A}
bir \lc{Prop}'tur ve onu \Lean{}'de ispatlayabiliriz. Bir totolojiyi
ispatlamanın en doğrudan yolu, değerleme üzerinde durum ayırmaktır;
bu, doğruluk tablosunu satır satır kontrol etmenin ta kendisidir:

\leandosya[Totolojiler]{Bolum10}{totoloji_ornek}

\begin{figure}[htbp]
\centering
\begin{tikzpicture}
  \draw[black!30, thick, fill=white, rounded corners=8pt] (0,0) rectangle (8.6,4.0);
  \node[anchor=north west, font=\baslikfont\small\bfseries, text=black!70] at (0.15,3.95) {bütün değerlemeler};
  \draw[ikincil, very thick, fill=ikincil!10] (5.0,2.0) ellipse (2.8cm and 1.55cm);
  \node[font=\baslikfont\small\bfseries, text=ikincil] at (6.9,3.15) {$A$ doğru};
  \draw[anarenk, very thick, fill=anarenk!15] (4.6,1.85) ellipse (1.55cm and 0.95cm);
  \node[font=\baslikfont\small\bfseries, text=anarenk] at (4.6,1.85) {$\Gamma$ doğru};
  \node[font=\small, anchor=west, align=left] at (9.0,2.4)
       {$\Gamma \models A$:\\$\Gamma$'yı doğrulayan\\her değerleme\\$A$'yı da doğrular};
\end{tikzpicture}
\caption{Mantıksal sonuç: $\Gamma$'nın modelleri, $A$'nın modellerinin
içindedir.}
\label{fig:sonuc}
\end{figure}

\noindent Mantıksal sonuç ilişkisinin bir örneği, modus ponens
kuralının anlambilimsel karşılığıdır:

\leandosya[Anlambilimsel modus ponens]{Bolum10}{sonuc_ornek}

\noindent Bir varsayım eklemek bir mantıksal sonucu bozmaz. Bu özelliğe
\emph{tekdüzelik} (İng.\ \emph{monotonicity}) denir:

\leandosya[Tekdüzelik]{Bolum10}{zayiflatma}

% ---------------------------------------------------------------------
\section{Bir karar yöntemi}\label{sec:karar-yontemi}

Bir formülün totoloji olup olmadığını \emph{hesaplayarak} anlayabiliriz:
İçinde geçen atomlara verilebilecek bütün değerleri tek tek deneriz.
$n$ atom için $2^n$ değerleme vardır:

\leandosya[Doğruluk tablosu yöntemi]{Bolum10}{karar}

\begin{figure}[htbp]
\centering
\begin{tikzpicture}[
    d/.style={circle, draw=anarenk, fill=tanimbg, thick, minimum size=0.55cm, font=\scriptsize},
    y/.style={rectangle, rounded corners=2pt, draw=#1, fill=#1!12, thick, font=\scriptsize, minimum width=1.0cm},
    level distance=1.1cm,
    level 1/.style={sibling distance=4.2cm},
    level 2/.style={sibling distance=2.1cm},
    edge from parent/.style={draw, black!45, thick},
  ]
  \node[d] {$p$}
    child {node[d] {$q$}
      child {node[y=bittiframe] {$\mathsf{T}$} edge from parent node[left, font=\scriptsize]{1}}
      child {node[y=bittiframe] {$\mathsf{T}$} edge from parent node[right, font=\scriptsize]{0}}
      edge from parent node[left=2pt, font=\scriptsize]{$p = 1$}}
    child {node[d] {$q$}
      child {node[y=bittiframe] {$\mathsf{T}$} edge from parent node[left, font=\scriptsize]{1}}
      child {node[y=bittiframe] {$\mathsf{T}$} edge from parent node[right, font=\scriptsize]{0}}
      edge from parent node[right=2pt, font=\scriptsize]{$p = 0$}};
  \node[font=\small, anchor=west, align=left] at (4.6,-1.1)
       {$n$ atom $\Rightarrow$ $2^n$ yaprak\\[2pt]
        {\scriptsize 10 atom: 1024 değerleme}\\
        {\scriptsize 50 atom: $\approx 10^{15}$ değerleme}};
\end{tikzpicture}
\caption{Doğruluk tablosu yöntemi bütün değerlemeleri dener. Peirce
yasası için 4 yaprağın hepsi ``doğru''dur. Yöntemin maliyeti atom
sayısıyla üstel olarak büyür.}
\label{fig:karar-agaci}
\end{figure}

\noindent \lc{totolojiMi} fonksiyonu yalnızca \lc{0, 1, …, n} atomlarını
dener; sonsuz sayıdaki öteki atomların değerlerini görmezden gelir. Bunun
neden yeterli olduğunu bir sonraki kesimde ispatlayacağız: Bir formülün
değeri yalnızca içinde geçen atomlara bağlıdır.

\begin{dikkat}[Dikkat: \#eval bir ispat değildir]
\lc{\textcolor{leancmd}{\bfseries\#eval} totolojiMi peirce} \lc{true}
sonucunu verir; ama bu, \lc{Totoloji peirce} önermesinin bir \emph{ispatı}
değildir. \lc{totolojiMi}'nin kendisi bir programdır ve hatalı olabilir.
Programı bir ispata dönüştürmek için onun \emph{doğruluğunu} ispatlamak
gerekir: \lc{totolojiMi A = true → Totoloji A}. Bu ispat, kendi
ispatlarını yapan karar yöntemlerinin (\emph{yansıma} tekniğinin)
temelidir; \ciltiki{6}'da tamamlayacağız.
\end{dikkat}

\begin{arastirmanotu}[Araştırma Notu: SAT ve P = NP]
Bir formülün \emph{sağlanabilir} olup olmadığını belirleme problemi
(SAT), Stephen Cook'un 1971'de gösterdiği gibi \emph{NP-tam} ilk problemdir
\citep{cook1971}. Doğruluk tablosu yöntemi üstel zaman alır; bundan
temelde daha iyi bir algoritmanın var olup olmadığı, bilgisayar biliminin
en ünlü açık problemi olan P = NP sorusuna denktir. Yine de modern
SAT çözücüleri, milyonlarca değişkenli endüstriyel formülleri
pratikte çözebilmektedir; bu konuya \ciltiki{6}'da döneceğiz.
\end{arastirmanotu}

% ---------------------------------------------------------------------
\section{İki önemli önsav}\label{sec:yerellik}

\subsection{Yerellik}

\begin{teorem}{Yerellik}{yerellik}
İki değerleme bir $A$ formülünde geçen bütün atomlarda aynı değeri
veriyorsa, $A$'ya da aynı değeri verir.
\end{teorem}

\leandosya[Yerellik]{Bolum10}{yerellik}

\noindent İspat, formülün yapısı üzerinde tümevarımdır. Her ikili bağlaç
durumunda, alt formüllerin atomları ana formülün atomlarının bir
alt kümesi olduğu için tümevarım hipotezleri uygulanabilir. Üç ikili
bağlacın durumları neredeyse aynıdır; bu tekrar, \ciltiki{7}'de
kendi taktiklerimizi yazmamız için iyi bir neden olacak.

\subsection{Yerine koyma}

Bir totolojide atomların yerine \emph{herhangi} formüller koyarsak yine
bir totoloji elde ederiz. Örneğin $p \oise p$ bir totoloji olduğuna göre
$(q \ove r) \oise (q \ove r)$ de bir totolojidir.

\leandosya[Yerine koyma]{Bolum10}{yerine_koy}

\noindent \lc{deger\_yerineKoy} önsavı, yerine koymanın anlamını tam
olarak söyler: $\sigma$ ile yerine koyduktan sonra $v$ altında
değerlendirmek, ``her atomu $\sigma$'daki formülün $v$ altındaki
değerine götüren'' yeni bir değerleme altında değerlendirmekle aynıdır.

% ---------------------------------------------------------------------
\section{Doğruluk tabloları klasiktir}\label{sec:klasik-anlambilim}

\lc{\textcolor{leancmd}{\bfseries\#eval} totolojiMi ucuncuHal} ve
\lc{totolojiMi peirce} sonuçları \lc{true}'dur. Oysa Bölüm~4'te bu iki
formülün \emph{yapıcı} olarak ispatlanamadığını görmüştük. Bu bir çelişki
değildir; bize anlambilimimiz hakkında önemli bir şey söyler:

\begin{onemlifikir}[Önemli Fikir: Her mantığın kendi anlambilimi vardır]
Doğruluk değerlerinin yalnızca iki eleman (\lc{true}, \lc{false})
içeren bir kümeden geldiği anlambilim, \emph{klasik} mantığın
anlambilimidir. Bu anlambilimde her formül ya doğru ya yanlıştır;
dolayısıyla $p \oveya {\sim}p$ her zaman doğrudur. Yapıcı mantığın
anlambilimi ise Bölüm~4'te önizlediğimiz Kripke modelleridir;
\ciltiki{2}'de onları da \Lean{}'de tanımlayacağız. Geçerlilik ve
tamlık teoremleri her zaman \emph{bir ispat sistemi} ile \emph{bir
anlambilim} arasındaki ilişki hakkındadır.
\end{onemlifikir}

\noindent Bu yüzden Bölüm~11 ve 12'de inceleyeceğimiz ispat sistemi,
klasik doğal çıkarım olacaktır: Bölüm~5'teki kural kataloğu ve RAA kuralı.

% ---------------------------------------------------------------------
\begin{arastirmanotu}
\textbf{Formalizasyonda sözdizimi.} Önermeler mantığında değişken
bağlama olmadığı için sözdizimi basittir. Yüklemler mantığında ise
niceleyicilerin bağladığı değişkenleri temsil etmek zordur; Bölüm~6'daki
değişken yakalama sorunu formalizasyonda ciddi bir mühendislik
problemine dönüşür. Yaygın çözümler de Bruijn indeksleri, ``yerel
adsız'' (\emph{locally nameless}) temsil ve nominal tekniklerdir.
Lean~4'teki Foundation projesi, Gödel teoremlerini formalize ederken
yerel adsız temsili kullanır \citep{saitou2026, foundation}.

\textbf{Derin ve sığ gömme.} Bu bölümde nesne dilini bir veri tipi
olarak tanımladık; buna \emph{derin gömme} (İng.\ \emph{deep embedding})
denir. Önceki bölümlerde ise mantığı doğrudan \Lean{}'in \lc{Prop}'u
içinde kullanıyorduk; buna \emph{sığ gömme} (\emph{shallow embedding})
denir. Derin gömme, mantık \emph{hakkında} teorem ispatlamamızı sağlar;
sığ gömme ise mantık \emph{içinde} ispat yapmayı kolaylaştırır.

\textbf{İleri okuma.} Önermeler mantığının anlambilimi ve metamantığı
için \citet{vandalen2013}; formalizasyon odaklı bir sunum için
\citet{harrison2009}.
\end{arastirmanotu}

\begin{bolumozeti}
\begin{itemize}
  \item Nesne dili incelenen mantıktır; üst dil, onun hakkında konuşulan
        dildir. \Lean{}'de üst dil \Lean{}'in kendisi, nesne dili ise
        \lc{Formul} veri tipidir.
  \item Formüller tümevarımsal bir tiptir; üzerlerinde özyineleme ve
        yapısal tümevarım yapılır.
  \item Bir değerleme atomlara doğruluk değeri verir; \lc{deger},
        Tarski'nin doğruluk tanımıdır.
  \item Totoloji, sağlanabilirlik ve mantıksal sonuç ($\Gamma \models A$)
        üst dilde tanımlanan önermelerdir.
  \item Doğruluk tablosu yöntemi totolojileri hesaplayarak bulur, ama
        maliyeti üsteldir. \lc{\#eval} bir ispat değildir.
  \item Yerellik ve yerine koyma önsavları, yapısal tümevarımla ispatlanır.
  \item İki değerli anlambilim klasiktir; yapıcı mantığın anlambilimi
        Kripke modelleridir.
\end{itemize}
\end{bolumozeti}

\begin{alistirmalar}
\alistirma{1} ${\sim}A$'nın boyutu, $A$'nın boyutundan bir fazladır.
(\lc{a10\_1})

\alistirma{2} $p \ove q \oise q \ove p$ bir totolojidir. (\lc{a10\_2})

\alistirma{2} $p \ove q \models p$. (\lc{a10\_3})

\alistirma{2} $\Gamma \models A$ ise $\Gamma \models A \oveya B$.
(\lc{a10\_4})

\alistirma{2} Bir totoloji her $\Gamma$'nın mantıksal sonucudur.
(\lc{a10\_5})

\alistirma{3} Anlambilimsel modus ponens: $\Gamma \models A \oise B$ ve
$\Gamma \models A$ ise $\Gamma \models B$. (\lc{a10\_6})

\alistirma{3} \textbf{Anlambilimsel çıkarım teoremi}: $A, \Gamma \models B$
ise $\Gamma \models A \oise B$. Bu teorem, Bölüm~5'teki $\to$G kuralının
anlambilimsel karşılığıdır. (\lc{a10\_7})

\alistirma{2} $\bot$ sağlanabilir değildir. (\lc{a10\_8})

\alistirma{1} \lc{\#eval} ile aşağıdaki formüllerin totoloji olup
olmadığını denetleyin: $(p \oise q) \oise ({\sim}q \oise {\sim}p)$,
$({\sim}q \oise {\sim}p) \oise (p \oise q)$, $(p \oise q) \oveya (q \oise p)$.
Hangileri Bölüm~4'e göre yapıcı olarak ispatlanamaz?

\alistirma{2} (Kâğıt-kalem) $\{\oise, \bot\}$ bağlaçlarının tek başına
yeterli olduğunu gösterin: $A \ove B$ ve $A \oveya B$'yi, yalnızca $\oise$
ve $\bot$ kullanarak, bütün değerlemelerde aynı değeri alan formüllerle
ifade edin.

\alistirma{3} (Proje) \lc{totolojiMi A = true → Totoloji A} önermesini
ispatlayın. İpucu: Önce \lc{tumDegerlemeler n}'nin, \lc{0, …, n-1}
atomlarında her olası değer birleşimini içerdiğini gösterin; sonra
yerellik önsavını kullanın.
\end{alistirmalar}
